\documentclass[11pt,a4paper,fontset=none,scheme=plain]{ctexart}


\usepackage[margin=25mm]{geometry}
\usepackage{amsmath,amssymb}
\usepackage{mathptmx}
\setmainfont{Times New Roman}
\setCJKmainfont[ItalicFont={Noto Serif CJK SC},BoldItalicFont={Noto Serif CJK SC Bold}]{Noto Serif CJK SC}
\setCJKsansfont{Noto Sans CJK SC}
\setCJKmonofont{Noto Sans Mono CJK SC}
\renewcommand{\abstractname}{摘要}
\renewcommand{\refname}{参考文献}
\renewcommand{\figurename}{图}
\renewcommand{\tablename}{表}
\usepackage{booktabs}
\usepackage{graphicx}
\usepackage{placeins}
\usepackage[round,authoryear]{natbib}

\usepackage{hyperref}
\hypersetup{colorlinks=true,linkcolor=black,citecolor=black,urlcolor=blue,
  pdftitle={多级差分电阻电容网络的端接约束与动态比值误差}}
\providecommand{\doi}[1]{\href{https://doi.org/#1}{\nolinkurl{doi:#1}}}
\setlength{\emergencystretch}{1.5em}
\clubpenalty=10000
\widowpenalty=10000
\displaywidowpenalty=10000
\raggedbottom
\title{多级差分电阻电容网络的端接约束与动态比值误差}
\author{}
\date{}
\begin{document}
\maketitle
\begin{abstract}
本文推导任意有限级数对地差分电阻电容梯形网络的端接相关共模抵消条件及动态比值误差。两个递推多项式统一表示开路电压传递函数与输出阻抗，逐级恢复检验给出正值元件可实现的充要条件。对于固定的 $N$ 级参考支路，开路匹配保留 $N$ 个局部自由度，指定一个正值电阻端接后保留 $N-1$ 个，在两个不同已知端接下完整传递函数均相同则要求对应元件相等。多项式核进一步将维数计算扩展到有理负载。对于缓冲后的比值测量，直流归一化传递函数相等时，常值电阻读出可精确抑制激励变化；变化电阻的读数则是受电流加权的历史平均。矩展开将电阻跟踪项与通道时延失配导致的有符号电流变化率项分开。覆盖一至十二阶的84例交流仿真通过独立节点方程核查；另有八例三级、四级瞬态验证，其中含两例步长细化，核验了历史权重和失配项。负载改变后的二阶支路对还产生共模阶跃泄漏脉冲，其单位阶跃峰值为 $9.197601$~mV/V，渐近时间常数为 $20~\mu\mathrm s$。这些结果将可实现元件族、端接和动态保真联系在同一网络表示中。
\end{abstract}
\noindent\textbf{关键词：}差分 RC 滤波器；共模抑制；电阻端接；输出阻抗；元件恢复。


\section{引言}

差分滤波器分别处理两个输入电压，再形成输出差值。两条通道的处理特性不一致时，共模输入会转化为差分输出。\citet{casas1996basics} 总结了两种共模抑制方法：匹配非耦合通道的传递函数，以及通过浮置公共节点耦合单端滤波器。本文研究采用对地支路的梯形网络，属于前一种配置。

需要匹配的电路量取决于测量频率和所接负载。\citet{hornero2021cm} 表明，元件失配可能使共模响应出现增益峰值和正相移。无源梯形网络的负载还会经输出阻抗取电流。因此，根据开路电压传递函数选择元件组合后，还需结合输出阻抗预测带载响应。

缓冲后的电阻传感接口中，样品电压与 $I(t)R(t)$ 成正比，参考电压与 $I(t)R_s$ 成正比。电阻在观测窗口内变化时，先滤波再取比会引入历史权重；通道响应失配时，电流变化还会影响常值电阻读数。因此，匹配规范需要同时确定线性网络的端接条件，以及动态比值实际估计的量。

本文计算指定端接下多级对地梯形网络可用于匹配的元件自由度，并将匹配条件与动态比值测量联系起来。利用经典梯形提取，将正值元件实现表示为系数坐标，进而分类讨论端接约束下的原始响应匹配与校准后匹配。精确构造和理想元件电路仿真核查负载依赖性；卷积分析区分常值电阻下的激励抑制、电阻历史权重，以及通道时延失配的一阶影响。

\begin{figure}[htbp]
\centering
\includegraphics[width=\linewidth]{figures/Fig_11_zh.pdf}
\caption{缓冲后比值模型的信号路径。高亮块为系数方程分析的无源网络，负载与探头导纳在其输出平面计入；传感电阻改变上游源幅值。}
\label{fig:system-plane}
\end{figure}


\section{相关研究与设计问题}
\label{sec:related-work}

\subsection{差分传递函数与阻抗平衡}
差分测量既要求保留有效的差模响应，也要求减小共模转差模响应。Pall\`as-Areny 和 Webster 建立了具有有限差模增益与共模增益的差分放大器模型，并通过实验验证所预测的幅值和相位\citep{pallasareny1991cmrr}。其配套分析以四个传递函数描述级联差分级，并明确假设后级不会对前级产生负载作用\citep{pallasareny1991cascaded}。对于直接连接接收器的无源梯形网络，接收器阻抗会进入支路传递函数，因此需要从建立匹配方程时就纳入负载。

复增益保留了比较共模抵消条件所需的信息。Koide 利用差模增益与共模增益的复数比，解释了有意引入不平衡如何补偿有源器件的共模误差\citep{koide1997redef}。Gasulla 等指出，描述全差分电路的瞬态行为和稳定性需要同时考虑共模增益与共模转差模增益\citep{gasulla2002cmresp}。Hornero 等进一步分析了支路失配引起的频率相关共模响应，包括带通行为和相移\citep{hornero2021cm}。这些复响应分析说明，需要分别报告各增益随频率变化的复响应；仅凭抑制比的幅值无法确定哪些支路系数必须相等。

\subsection{拓扑、端接与元件匹配}
Casas 和 Pall\`as-Areny 分析了非耦合与耦合模拟差分滤波器，并给出了高阶扩展\citep{casas1996basics}。非耦合支路通过传递函数匹配实现抵消；改变互联方式还可利用电路拓扑实现抵消。其中，浮置公共节点结构为基于拓扑的共模抑制提供了重要参照。图示二阶实例由两个对地梯形网络组成，其输出端可连接等值电阻负载，并在两个输出之间接入跨接电容。因此，元件条件需要针对这些具体连接关系推导。

阻抗平衡也决定实际平衡接口的响应。Whitlock 和 Floru 展示了信号源与接收器的共模阻抗不平衡如何相互作用，并介绍了一种具有高共模输入阻抗的电路\citep{whitlock2004balancedinput}。对于本文梯形网络，需要回答的问题是：两条开路电压传递函数相同的支路，在接入负载后是否仍然相同？开路传递函数确定电压比，输出阻抗确定所接负载如何改变这一电压比。保留这两个函数，便可显式描述端接的影响，直接比较开路匹配与指定有限负载下的匹配。

\subsection{网络实现与容差分析}
经典综合的目标是由给定传递函数构造一个网络。Guillemin 建立了连分式梯形综合，并给出端接后的传递结构 $Y_{12}/(1+Y_{22})$\citep{guillemin1949rcsynthesis}。Fialkow 和 Gerst 刻画了梯形网络与一般双端口 RC 网络的传递函数及增益常数的可实现范围\citep{fialkow1951rcladder,fialkow1952rcnetworks}；其连续实现族构造还表明，同一传递函数允许许多正值实现。梯形网络的最大可实现增益由 Paige 和 Kuh 给出，Wang 等借助 Ali 和 Kim 的同余变换证明该范围内任意指定增益均可实现\citep{paige1960maximum,wang2017rcladder,ali1975continuous}。端接在这一理论中以影响可实现增益范围的方式进入\citep{fialkow1951rcladder,fialkow1952rcnetworks}。

本文提出的匹配问题与此不同：本文计算在指定端接和测量口径下带载传递函数等于固定参考支路的正值元件族。传输矩阵与连分式表示是完成该计算的已有工具\citep{guillemin1949rcsynthesis,fialkow1951rcladder,lee1978chain}，本文用其表达已知端接导纳施加的约束。所得计数区分包含绝对增益的匹配与直流校准后的匹配，不属于新的可实现性定理。

元件实现需要与灵敏度一起考虑。Moschytz 讨论了有源滤波器设计中灵敏度、元件值跨度与元件容差之间的权衡\citep{moschytz0radeffs}。Director 和 Rohrer 建立了用于计算网络灵敏度的广义伴随网络框架\citep{director1969the}。Endres 和 Kirkpatrick 研究了差分环路滤波器失配如何在快速建立的锁相环中引入额外瞬态\citep{endres0ensitivi}。这些灵敏度与瞬态研究支持将精确标称抵消条件、元件扰动影响和建立时间规范分别考察；各问题均需要明确的响应函数与评价准则。

\subsection{本文工作的范围}
本文推导任意有限级数正值对地 RC 梯形网络带载传递函数相等的系数条件。经典梯形提取提供系数坐标，多项式核给出已知有理负载下的局部匹配维数。分析进一步将原始抵消条件与校准后的动态比值测量联系起来，解析构造通过原生交流与瞬态电路输出核验。

上述抵消条件针对确定性的线性传递函数。Serrano-Finetti 等表明，共模电子噪声分析除确定性共模抑制比外，还需要噪声源及其互相关模型\citep{serranofinetti2019cmnoise}。因此，评价噪声性能需要给出相关噪声源的位置与相关性。第~\ref{sec:noise-pairing} 节在指定温度和频带下计算电阻热噪声，并保留比值读出中的通道协方差项。

\section{网络模型与验证方法}

\subsection{信号通道与测量平面}
\label{sec:model-plane}

每条信号通道为 $N$ 级梯形网络，由正值串联电阻和正值对地电容组成，并由理想电压源驱动。图~\ref{fig:circuit} 展示二级实例，输出取末端节点电压差。传递函数按零初始条件定义。缓冲后的比值模型中，样品电阻先改变输入源幅度，滤波元件保持固定；输入前史确定相容的初始状态。

差模输入和共模输入定义为

\begin{equation}\label{eq:1}
v_d=v_p-v_n,\qquad v_c=\frac{v_p+v_n}{2}.
\end{equation}

在差模输入为零时，共模转差模增益为 \(G_{DC}\) = \(v_{d,\mathrm{out}}/v_{c,\mathrm{in}}\)；在共模输入为零时，差模增益为 \(G_{DD}\) = \(v_{d,\mathrm{out}}/v_{d,\mathrm{in}}\)。共模抑制比为 \(\lvert G_{DD}/G_{DC}\rvert\)。对于传递函数分别为 \(H_P\) 和 \(H_N\) 的独立通道，

\begin{equation}\label{eq:2}
G_{DC}=H_P-H_N,\qquad
G_{DD}=\frac{H_P+H_N}{2}.
\end{equation}

共模抵消要求 \(G_{DC}\) 作为频率的有理函数恒等于零。

\begin{figure}[htbp]
\centering
\includegraphics[width=0.86\linewidth]{figures/Fig_01_circuit.pdf}
\caption{两条对地二阶 RC 支路及输出端接。每条梯形支路含两个串联电阻和两个对地电容。可选电容 $C_d$ 跨接输出，跨接路径以蓝色点线表示。两个输入由对地理想电压源驱动。}
\label{fig:circuit}
\end{figure}


\subsection{系数表示与负载条件}
\label{sec:coefficients}

以 \(R_*\) 和 \(C_*\) 为基准归一化元件值，定义

\begin{equation}\label{eq:3}
r_i=R_i/R_*,\qquad q_i=C_i/C_*,\qquad x=sR_*C_*.
\end{equation}

对于二阶示例，支路的四个系数为

\begin{equation}\label{eq:4}
\begin{aligned}
a&=r_1q_1+(r_1+r_2)q_2,\\
b&=r_1r_2q_1q_2,\\
c&=r_1+r_2,\\
d&=r_1r_2q_1.
\end{aligned}
\end{equation}

由两个节点方程可得开路电压传递函数 \(H_0\) 和输出阻抗 \(Z_0\)：

\begin{equation}\label{eq:5}
H_0(x)=\frac{1}{1+ax+bx^2},\qquad
Z_0(x)=R_*\frac{c+dx}{1+ax+bx^2}.
\end{equation}

输出端与地之间接入电阻 \(R_L\) 时，其归一化电导为 \(g=R_*/R_L\)，带载传递函数为

\begin{equation}\label{eq:6}
H_g(x)=\frac{H_0(x)}{1+Z_0(x)/R_L}
      =\frac{1}{1+gc+(a+gd)x+bx^2}.
\end{equation}

每一种负载条件下，两条通道均连接相同的指定电阻 \(R_L\)。统计匹配支路 N 的可行元件组合时，参考支路 P 保持固定。每个电导条件下比较的对象均为完整传递函数；比较两个电导条件，即分别比较每条通道在这两种负载下随频率变化的传递函数。

\subsection{精确运算与数值验证}

采用精确有理数运算核查系数恒等式和示例元件组合。电路仿真使用 LTspice 26.1.1，取 \(R_*=1\,\mathrm{k\Omega}\)、\(C_*=10\,\mathrm{nF}\)。扫频范围为 10 Hz 至 10 MHz，每十倍频程 24 点，共 145 点。共模激励在两个输入端同时施加 1 V；差模激励分别施加 +0.5 V 和 -0.5 V。

37 个仿真算例覆盖一阶对照、二阶支路、不同元件匹配策略、有限源电阻、对称与非对称电阻端接，以及输出跨接电容。根据保存的网表独立建立节点方程，将所得复节点电压与仿真输出逐点比较。电阻负载分类的解析结论采用第~\ref{sec:model-plane}、\ref{sec:coefficients} 节所述的理想电压源和两侧等值电阻负载假设。


一般阶矩阵包括 $N=1,2,3,4,6,8,12$ 的84例原生交流仿真。独立节点核查采用实际元件值及仿真器规定的元件模型；其中，未显式给出串联电阻的电感在 LTspice 中默认串入1~m$\Omega$。另做两组成对诊断运行，分别显式指定零或1~m$\Omega$。瞬态验证采用第~\ref{sec:dynamic-verif} 节所列的波形与热身窗口，两个四级算例的细化检查将最大步长减半。

\FloatBarrier

\section{任意级数下的抵消条件}
\label{sec:general-order}

\subsection{梯形网络的递推表示}

考虑包含 $N\geq1$ 级的支路，每级均由正值串联电阻 $r_j$ 及其后的正值对地电容 $q_j$ 组成。元件归一化方法和频率变量沿用式~\eqref{eq:3}。图~\ref{fig:circuit} 中的二级电路属于这类网络。取电流从网络流向端接负载为正方向，第 $j$ 级的传输矩阵为
\begin{equation}\label{eq:general-matrix}
M_j(x)=\begin{bmatrix}1+r_jq_jx&r_j\\q_jx&1\end{bmatrix}.
\end{equation}
设 $(A_j,B_j)$ 为 $M_1\cdots M_j$ 的第一行。接入末级后，有
\begin{equation}\label{eq:general-forward}
B_j=r_jA_{j-1}+B_{j-1},\qquad
A_j=A_{j-1}+xq_jB_j,
\quad A_0=1,\ B_0=0.
\end{equation}
由式~\eqref{eq:general-forward} 归纳可得 $\deg A_N=N$、$\deg B_N=N-1$ 且 $A_N(0)=1$。两个多项式的最高次项系数均为正。各级矩阵的行列式均为 1，其乘积的行列式也为 1，因此第一行的两个多项式不存在非常数公因子。

由端口方程可得开路传递函数和输入源短路时的输出阻抗。在两条支路分别接入相同电阻负载后，即得到下述带载传递函数；该端接分解与经典综合中的结构一致\citep{guillemin1949rcsynthesis}：
\begin{equation}\label{eq:general-transfer}
H_0=\frac1{A_N},\qquad
\frac{Z_0}{R_*}=\frac{B_N}{A_N},\qquad
H_g=\frac1{A_N+gB_N}.
\end{equation}
当 $N=2$ 时，$A_2=1+ax+bx^2$、$B_2=c+dx$，即得到式~\eqref{eq:5}。保留 $B_N$ 可描述开路电压传递函数未包含的端接信息。

\subsection{经典梯形提取与系数坐标}

输入源短路后，$B_N/A_N$ 为从输出端看入的无量纲驱动点阻抗，可写成梯形连分式\citep{guillemin1949rcsynthesis,lee1978chain}
\[
z_j=\left[q_jx+\left(r_j+z_{j-1}\right)^{-1}\right]^{-1},
\qquad z_0=0,\quad z_N=B_N/A_N.
\]
下面的逐级恢复将已有连分式梯形综合\citep{guillemin1949rcsynthesis,lee1978chain} 特化为单电阻/单电容臂拓扑的 $(A_j,B_j)$ 坐标；一般臂梯形网络的递推传递计算也属经典方法\citep{fialkow1951rcladder}。本文利用该提取结果，计算指定端接约束下的可实现元件族维数。

末级元件值可由最高次项系数恢复。以 $\operatorname{lc}(p)$ 表示多项式 $p$ 的最高次项系数，定义
\begin{equation}\label{eq:general-inverse}
\begin{aligned}
q_j&=\frac{\operatorname{lc}(A_j)}{\operatorname{lc}(B_j)},&
A_{j-1}&=A_j-xq_jB_j,\\
r_j&=\frac{\operatorname{lc}(B_j)}{\operatorname{lc}(A_{j-1})},&
B_{j-1}&=B_j-r_jA_{j-1}.
\end{aligned}
\end{equation}
对于全正值元件梯形网络，上述运算恰好逆转式~\eqref{eq:general-forward} 的递推。每一步中，$A_{j-1}$ 的次数为 $j-1$；当 $j>1$ 时，$B_{j-1}$ 的次数为 $j-2$。最终得到 $(A_0,B_0)=(1,0)$。各比值均为正值，且唯一确定。

反之，从满足上述次数要求且 $A_N(0)=1$ 的实系数多项式出发，若每次相减均得到要求的次数、恢复出的每个元件值均为正，且最终得到 $(1,0)$，则逆序执行上述运算，即可由全正值元件梯形网络重构原多项式。因此，这些条件给出了该拓扑可实现性的有限步充要检验。正值或次数检查中任一项不满足，即可排除该候选系数对。

从 $2N$ 个正值元件到 $A_N$ 和 $B_N$ 的 $2N$ 个系数的正向映射为多项式映射。式~\eqref{eq:general-inverse} 给出的逆映射为有理映射，在每个正值元件实现处，其分母均非零。在足够小的邻域内，元件值仍严格为正，最高次项系数仍非零。两个映射在该邻域内光滑且互逆，因此这些系数可作为元件族的局部坐标。

\subsection{匹配自由度的数目}

固定一条参考支路。开路匹配固定了 $A_N$ 的 $N$ 个非常数项系数，在正值元件可实现域内保留 $N$ 个局部自由度。同一传递函数具有多个正值实现属经典结论\citep{fialkow1951rcladder,fialkow1952rcnetworks}；此处计数的是本拓扑内与固定参考支路匹配的元件族。统一阻抗缩放使 $B_N$ 乘以同一个标量因子，同时保持 $A_N$ 不变。当 $N\geq2$ 时，该缩放仅沿更大匹配族中的一个方向变化。

在指定的 $g>0$ 下，原始响应匹配 $H_{g,P}=H_{g,N}$ 固定 $D_g=A_N+gB_N$ 的全部系数。已有综合理论以限制可实现增益范围的方式处理端接\citep{fialkow1951rcladder,fialkow1952rcnetworks}；以下计数针对的则是与给定参考支路匹配的元件族。常数项通过 $D_g(0)=1+gB_N(0)$ 固定 $B_N(0)$，其余 $N$ 个系数固定 $A_N$ 与 $B_N$ 系数的组合。这 $N+1$ 个独立线性约束保留 $N-1$ 个局部自由度，$N=1$ 时为零。

直流校准后比较的对象为 $G_g=D_g(0)/D_g$。固定归一化分母产生 $N$ 个独立约束，保留 $N$ 个局部自由度。在系数坐标中，可以通过改变 $A_N$，独立改变归一化分母的最高次项及其余各项系数，故约束映射的秩为 $N$。此处使用的增益自由度在指定增益综合中已有对应结论：固定极点后，增益常数可在可实现区间内变化\citep{fialkow1951rcladder,paige1960maximum,wang2017rcladder}。在本文的匹配问题中，该自由度表现为直流校准后配对支路的歧义，并影响未经校准的差分输出。对于 $N=1$，正值元件族
\[
r>0,\qquad q=\frac{r+10}{11r}
\]
在 $g=0.1$ 时与 $(r,q)=(1,1)$ 具有相同归一化响应，但其原始直流增益随 $r$ 改变。

对于 $N=2$，支路 P = $(1,1,1,1)$ 与正值支路
\[
\left(2,2,\frac{217+\sqrt{6769}}{288},\frac{217-\sqrt{6769}}{480}\right)
\]
在 $g=0.1$ 时具有相同归一化响应。两者原始分母相差常数倍 $7/6$。写成 $D_N=(7/6)D_P$ 可得 $H_N=(6/7)H_P$，故原始共模转差模增益为 $G_{DC}=H_P-H_N=H_P/7$，差模增益为 $G_{DD}=(H_P+H_N)/2=(13/14)H_P$。两者之比与频率无关，原始 CMRR 在所有频率均为 $20\log_{10}|G_{DD}/G_{DC}|=20\log_{10}(13/2)=16.26$~dB。校准因此允许额外的匹配方向，该方向并不消除原始差分输出：读出所用的归一化响应一致，而原始共模转差模转换仍然存在。

\begin{table}[htbp]
\centering
\caption{固定正值 $N$ 级参考支路时，电阻端接下的局部匹配维数。原始匹配包含绝对增益，归一化匹配比较 $H_g/H_g(0)$。}
\label{tab:matching-conventions}
\begin{tabular}{lll}
\toprule
端接信息 & 原始响应 & 归一化响应 \\
\midrule
开路 & $N$ & $N$ \\
一个指定 $g>0$ & $N-1$ & $N$ \\
两个不同指定 $g_1,g_2>0$ & $0$ & $0$ \\
\bottomrule
\end{tabular}
\end{table}

对于两个不同的正值电导，令 $U_i=1/H_{g_i}$。两个完整传递函数可确定
\begin{equation}\label{eq:general-two-load}
B_N=\frac{U_2-U_1}{g_2-g_1},\qquad
A_N=\frac{g_2U_1-g_1U_2}{g_2-g_1}.
\end{equation}
因此，若两条支路在两种负载下均匹配，则其 $A_N$ 和 $B_N$ 分别相同。由经典梯形提取的逐步唯一性可知，对应元件值必须相同\citep{guillemin1949rcsynthesis}。该论证适用于上述梯形网络类别中的任意有限 $N$，所用信息为两种负载下包含绝对增益的完整复传递函数。

两个归一化电阻端接响应也可精确确定元件。对候选支路，写成 $D_i(g_k)=c_kD_{\rm ref}(g_k)$。$x^N$ 的系数不随 $g_k$ 改变且非零，因此 $c_1=c_2=c$。两负载方程相减得 $B_i=cB_{\rm ref}$，进而得 $A_i=cA_{\rm ref}$。由 $A_i(0)=A_{\rm ref}(0)=1$ 可知 $c=1$。该证明使用精确完整响应，尚未规定含噪数据的稳定估计方法。

\subsection{负载类型对匹配族的影响}

已知的有理函数负载可用无量纲导纳 $g(x)=P(x)/Q(x)$ 表示，其中 $P$ 和 $Q$ 为非零且互素的实系数多项式。假设 $QA_N+PB_N$ 不恒为零，使带载响应有定义。两条支路的传递函数相等可得
\begin{equation}\label{eq:rational-kernel}
H_g=\frac{Q}{QA_N+PB_N},\qquad
Q\Delta A+P\Delta B=0.
\end{equation}
其中 $\Delta$ 表示支路 P 减去支路 N。由互素性可知，存在多项式 $S$，使 $\Delta A=PS$ 且 $\Delta B=-QS$。次数上限和 $\Delta A(0)=0$ 对 $S$ 的约束为
\begin{equation}\label{eq:rational-dimension}
\begin{aligned}
m&=\min\{N-\deg P,\ N-1-\deg Q\},\\
d_g&=\max\{0,m+1-\epsilon\},\qquad
\epsilon=\begin{cases}1,&P(0)\ne0,\\0,&P(0)=0.\end{cases}
\end{aligned}
\end{equation}
当 $m<0$ 时，仅允许 $S=0$。否则，$S$ 有 $m+1$ 个系数；若 $P(0)$ 非零，则其常数项系数必须为零。该线性族与正值元件可实现开域相交后，局部维数仍为 $d_g$。

表~\ref{tab:load-dimensions} 给出将式~\eqref{eq:rational-dimension} 用于常见负载的结果。开路时 $P$ 恒为零，单独列出。纯电容负载改变端接电容，保留 $N$ 个自由度；电阻负载还约束零频增益，将自由度减少至 $N-1$。
\begin{table}[htbp]
\centering
\caption{在一种指定负载下、固定参考支路时的原始响应匹配族局部维数，包含绝对增益约束。表中各负载参数均为正值。直流归一化匹配另见表~\ref{tab:matching-conventions}。}
\label{tab:load-dimensions}
\begin{tabular}{lll}
\toprule
负载 & 无量纲导纳 & 维数 \\
\midrule
开路 & $0$ & $N$ \\
电阻 & $g_0$ & $N-1$ \\
电容 & $\alpha x$ & $N$ \\
电感 & $\beta/x$ & $\max(N-2,0)$ \\
串联 RC & $\alpha x/(1+\gamma x)$ & $N-1$ \\
并联 RC & $g_0+\alpha x$ & $N-1$ \\
\bottomrule
\end{tabular}
\end{table}

在原始响应匹配口径下，两个不同的已知有理导纳也可确定元件组合。令 $g_k=P_k/Q_k$，分别将 $Q_k\Delta A+P_k\Delta B=0$ 乘以另一负载的分母后相减，得到 $(P_1Q_2-P_2Q_1)\Delta B=0$。负载函数不恒等意味着 $P_1Q_2-P_2Q_1\not\equiv0$，因此 $\Delta B=0$，进而 $\Delta A=0$。已知非零多项式 $Q_k$ 可在有理函数恒等式中约去，孤立零点处由延拓处理。表~\ref{tab:load-dimensions} 中的负载均为时不变有理导纳；周期开关输入需要第~\ref{sec:dynamic} 节另行引入的采样模型。

\subsection{负载变化与恢复对测量误差的敏感性}

对于独立支路，令 $D_i(g)=A_i+gB_i$。共模转差模增益的精确表达式为
\begin{equation}\label{eq:general-leakage}
G_{DC}=-\frac{\Delta A+g\Delta B}{D_P(g)D_N(g)}.
\end{equation}
若两条支路在 $g_0>0$ 下匹配，则 $\Delta A=-g_0\Delta B$ 且 $\Delta B(0)=0$。当两侧同时换用另一种相同电阻负载时，
\begin{equation}\label{eq:general-load-shift}
G_{DC}(g)=-\frac{(g-g_0)\Delta B}{D_P(g)D_N(g)}.
\end{equation}
零频泄漏仍为零，其余随频率变化的失配由多项式 $\Delta B$ 决定。二阶网络的负载变化公式对应 $\Delta B=x\Delta d$。

精确可辨识性与恢复对测量误差的敏感性需要分别检查。式~\eqref{eq:general-two-load} 以负载电导之差为分母，响应的小扰动则满足 $\delta(1/H)\simeq-\delta H/H^2$。两种负载过于接近或响应衰减过强，均会放大测量误差。因此，实验中的元件恢复除代数唯一性结论外，还需要明确频带和不确定性模型。
对无量纲系数映射 $(r_1,\ldots,r_N,q_1,\ldots,q_N)\mapsto(A_{N,1},\ldots,A_{N,N},B_{N,0},\ldots,B_{N,N-1})$，单位元件点未缩放Jacobian的2范数条件数在 $N=2,4,6,8,10,12$ 时分别为 $15.6$、$1.08\times10^3$、$9.30\times10^4$、$8.67\times10^6$、$8.40\times10^8$、$8.35\times10^{10}$。这些数值针对所列系数归一化，随坐标缩放改变，描述系数恢复的局部误差放大，并非特定响应拟合估计器的性能界。
\FloatBarrier

\section{抵消条件与元件值恢复}

\subsection{正值元件恢复}
图~\ref{fig:positive-region} 给出固定 $a$、$b$ 时的正值元件可实现域。

由四个系数可按下式确定元件值：

\begin{equation}\label{eq:7}
q_2=\frac{b}{d},\qquad
r_2=\frac{d^2}{ad-bc},\qquad
r_1=c-r_2,\qquad
q_1=\frac{d}{r_1r_2}.
\end{equation}

系数组合能够恢复为全正值元件的充要条件为

\begin{equation}\label{eq:8}
b>0,\qquad d>0,\qquad ad-bc>0,\qquad
c(ad-bc)>d^2.
\end{equation}

前两个不等式保证 \(q_2\) 为正，第三个保证 \(r_2\) 为正，第四个保证 \(r_1\) 为正，进而保证恢复出的 \(q_1\) 为正。将式~\eqref{eq:7} 代入式~\eqref{eq:4} 可还原原系数。反之，全正值元件满足式~\eqref{eq:8}，且有恒等式 \(ad-bc=r_1^2r_2q_1^2\)。因此，该映射在全正值元件域内为一一映射。

\begin{figure}[htbp]
\centering
\includegraphics[width=0.67\linewidth]{figures/Fig_02_positive_region.pdf}
\caption{开路系数固定为 $a=3$、$b=1$ 时的正值元件可实现区域。阴影内部满足 $d>0$ 及 $c(3d-c)>d^2$，边界线不属于该区域。P、Q 是同一开路传递函数的两组不同实现。坐标为归一化系数。}
\label{fig:positive-region}
\end{figure}


\subsection{开路共模抵消}

输出开路时，式~\eqref{eq:2} 中的 \(G_{DC}\) 恒为零的充要条件为

\begin{equation}\label{eq:9}
a_P=a_N,\qquad b_P=b_N.
\end{equation}

固定 a 和 b 后，c 和 d 仍可在式~\eqref{eq:8} 的可行域内变化。因此，等开路传递函数对应一个双参数元件族。统一阻抗缩放 \(r_i\rightarrow kr_i\)、\(q_i\rightarrow q_i/k\) 保持 a 和 b 不变，是该元件族中的一种构造。

另一组构造展示了统一缩放以外的自由度。按元件顺序 \((r_1,r_2,q_1,q_2)\)，支路 P = (1,1,1,1) 和 Q = (1,4,1/2,1/2) 均得到 a = 3、b = 1。Q/P 的两个电阻比分别为 1 和 4，无法由同一个阻抗缩放因子产生。P 和 Q 的输出阻抗系数 \((c,d)\) 分别为 (2,1) 和 (5,2)。

\subsection{单一正值负载电导下的共模抵消}

对于指定的 \(g>0\)，式~\eqref{eq:6} 相等要求

\begin{equation}\label{eq:10}
c_P=c_N,\qquad b_P=b_N,\qquad
 a_P+gd_P=a_N+gd_N.
\end{equation}

由于分子均为常数 1，传递函数相等与分母各项系数相等互为充要条件。式~\eqref{eq:10} 的三个条件在全正值元件域内保留一个自由系数。选定 d 后，可由指定的 a + gd 确定 a，再按式~\eqref{eq:7} 恢复元件值。

支路 P 和 T 给出了该元件族的实例。支路 T = (18/19,20/19,361/180,1/2) 对应 \((a,b,c,d)\) = (2.9,1,2,2)。当 g = 0.1 时，P 和 T 的分母均为 \(1.2+3.1x+x^2\)；当 g = 0.2 时，P 和 T 的分母分别变为 \(1.4+3.2x+x^2\) 和 \(1.4+3.3x+x^2\)。负载改变后，原先在第一种端接下匹配的两条响应发生分离。

\subsection{两个不同正值负载电导下的共模抵消}

设两个支路在 \(g_1>0\)、\(g_2>0\) 且 \(g_1\ne g_2\) 时均具有相同的完整传递函数。比较分母常数项和二次项可得 \(c_P\) = \(c_N\)、\(b_P\) = \(b_N\)；比较一次项可得

\begin{equation}\label{eq:11}
\Delta a+g_1\Delta d=0,\qquad
\Delta a+g_2\Delta d=0,
\end{equation}

其中 \(\Delta\) 表示两个支路对应系数之差。两式相减得到 \(\Delta d=0\)，进而有 \(\Delta a=0\)。至此，四个系数均相等，由式~\eqref{eq:7} 的逆映射可知，对应元件值必须相等。

对于固定参考支路，双负载要求将带有一个自由系数的可行匹配族收缩为唯一的元件组合。对应元件相同也直接满足该要求，从而建立了所定义梯形网络类别中的必要性与充分性。


\subsection{负载变化与元件偏差引起的泄漏}

以 $D_i(g)=1+gc_i+(a_i+gd_i)x+b_ix^2$ 表示支路 $i$ 的传递函数分母。两个支路传递函数之差的精确表达式为
\begin{equation}\label{eq:load-leakage}
G_{DC}=\frac{D_N(g)-D_P(g)}{D_P(g)D_N(g)}.
\end{equation}
对于在正值电导 $g_0>0$ 处匹配的支路，定义 $\Delta d=d_P-d_N$。由式~\eqref{eq:10} 可得
\begin{equation}\label{eq:load-shift}
G_{DC}(g)=-\frac{(g-g_0)\Delta d\,x}{D_P(g)D_N(g)}.
\end{equation}
分子明确给出了负载变化量与剩余输出阻抗系数差的共同作用。对于 $d$ 不相等的支路对，两侧负载同时改变后，直流泄漏仍为零，但非零频率处会出现随频率变化的泄漏。因此，仅匹配直流输出无法消除单一负载条件下的元件不唯一性。

对于匹配支路附近的小系数偏差，设其传递函数为 $H_g$，则一阶泄漏为
\begin{equation}\label{eq:sensitivity}
G_{DC}\simeq-H_g^2\left[g\Delta c+x(\Delta a+g\Delta d)+x^2\Delta b\right].
\end{equation}
其中，各差值均取支路 P 减去支路 N；在支路分母非零的频点，忽略项为系数偏差的二阶项。三个频率相关权重分别对应分母常数项、一次项和二次项的贡献，因此可直接针对带载系数组合分配容差。

即使两条支路的元件值完全相同，两侧外部负载不等也会产生共模转差模。两侧负载电导分别为 $g_P$ 和 $g_N$ 时，
\begin{equation}\label{eq:load-asymmetry}
G_{DC}=-\frac{(g_P-g_N)(c+dx)}{D(g_P)D(g_N)}.
\end{equation}
式~\eqref{eq:load-asymmetry} 将负载平衡纳入共模抵消规范。相同元件在相同端接下实现抵消；两侧实际接入的负载导纳也必须相等。

\FloatBarrier

\section{动态比值误差与共模瞬态泄漏}
\label{sec:dynamic}

第~\ref{sec:model-plane} 节的缓冲前端将 $v_P=IR$ 和 $v_N=IR_s$ 分别送入两条独立的全正值元件 RC 梯形网络，端接负载为恒定的非负电阻负载。电阻变化改变信号源幅度，各梯形网络仍为稳定的线性时不变系统。下文的卷积包含完整的输入历史，或采用与其等效且相容的初始状态，并要求参考通道输出非零。记 $H_j=1/D_j$、$G_j=H_j/H_j(0)=D_j(0)/D_j$，以 $g_j$ 表示 $G_j$ 的冲激响应。测得输出电压 $u_P,u_N$ 后，经直流增益校准得到
\begin{equation}\label{eq:ratio-readout}
\hat R=R_s\frac{H_N(0)}{H_P(0)}\frac{u_P}{u_N}
=\frac{(g_P\ast IR)(t)}{(g_N\ast I)(t)}.
\end{equation}

\subsection{匹配通道：精确的加权平均恒等式}
\label{sec:dynamic-matched}

当 $G_P=G_N=G$ 时，减去瞬时电阻可得
\begin{equation}\label{eq:weighted-average}
\hat R(t)-R(t)=\frac{\displaystyle\int_0^\infty g(u)I(t-u)[R(t-u)-R(t)]\,\mathrm du}
{(g\ast I)(t)}.
\end{equation}
对于正值激励，读数是电阻历史值的加权平均。当电阻恒为 $R_0$ 时，对于所有使分母非零的激励波形，读数均等于 $R_0$。当电阻变化时，激励仍参与确定权重。例如，取 $I=I_0e^{\alpha t}$、$R=R_0e^{\beta t}$，其中 $\alpha,\beta\ge0$，且完整历史均满足相应的指数关系，则有
\begin{equation}\label{eq:dynamic-exp}
\hat R(t)=R(t)\frac{G(\alpha+\beta)}{G(\alpha)}.
\end{equation}
即使通道完全相同，改变 $\alpha$ 通常仍会改变读数（图~\ref{fig:dynamic-weighting}）。

对于全正值元件的对地 RC 梯形网络，节点状态矩阵的非对角元素均非负，输入和输出向量也均非负。因此，其矩阵指数给出 $g(t)\ge0$。部分分式展开的留数符号交替，并不意味着冲激响应存在负值区段。于是，
\begin{equation}\label{eq:dynamic-moment}
M_1=\int_0^\infty u|g(u)|\,\mathrm du
=\int_0^\infty ug(u)\,\mathrm du=\tau,
\qquad \tau=T\frac{D_x'(0)}{D(0)},\quad T=R_*C_*.
\end{equation}
若在相关输入历史内满足 $|\dot R|\le L_R$ 和 $|I|\le I_{\max}$，则有
\begin{equation}\label{eq:lag-bound}
|\hat R-R|\le\frac{L_R I_{\max}\tau}{|(g\ast I)(t)|}.
\end{equation}
若在该段完整历史内始终有 $0<I_{\min}\le I\le I_{\max}$，则上界为 $L_R(I_{\max}/I_{\min})\tau$。若从有限时刻启动，则需要使用实际分母。

\subsection{失配通道：电流变化率的贡献}
\label{sec:dynamic-mismatch}

归一化响应核不同时，误差的精确分解为
\begin{equation}\label{eq:dynamic-decomposition}
\hat R-R=\frac{\displaystyle\int_0^\infty g_P(u)I(t-u)[R(t-u)-R(t)]\,\mathrm du}{(g_N\ast I)(t)}
+R\frac{((g_P-g_N)\ast I)(t)}{(g_N\ast I)(t)}.
\end{equation}
两项均采用参考通道的分母。当 $I,R$ 缓慢变化，且分母的绝对值具有大于零的下界时，展开各卷积可得
\begin{equation}\label{eq:mismatch-dynamic}
\hat R-R\simeq-\tau_P\dot R-R\frac{\dot I}{I}\Delta\tau,
\qquad \Delta\tau=\tau_P-\tau_N.
\end{equation}
其中 $\tau_j=T(A_{j,1}+gB_{j,1})/(1+gB_j(0))$。可用 $\mu_{2,j}=\int u^2g_j(u)\,\mathrm du$ 控制展开余项：对于二阶可导输入 $f$，
\begin{equation}\label{eq:dynamic-remainder}
|(g_j\ast f)-f+\tau_j\dot f|
\le\tfrac12\mu_{2,j}\sup_{v\le t}|f''(v)|.
\end{equation}
先分别将该界用于 $f=IR$ 和 $f=I$，再展开二者的比值。对于快速变化的边沿，需要使用精确表达式。参考通道较慢时，$\Delta\tau<0$；在电阻恒定、电流上升时，会产生正误差（图~\ref{fig:dynamic-slew}）。

\subsection{恒定电阻下的匹配判据}
\label{sec:constant-resistance-match}

校准后的读数对所有允许的电流历史均等于任意恒定电阻 $R_0$ 的充要条件为 $g_P=g_N$。充分性由式~\eqref{eq:ratio-readout} 得到。对于必要性，若每个探测电流均满足 $(g_P-g_N)\ast I=0$，则两个响应核必须相等。在正值恒定电流附近施加小扰动，已足以作为探测信号。等价地，
\begin{equation}\label{eq:dynamic-normalized-match}
G_P=G_N\quad\Longleftrightarrow\quad
\frac{D_P}{D_P(0)}=\frac{D_N}{D_N(0)}.
\end{equation}
未归一化分母相等，即 $D_P=D_N$，还要求两个直流增益相等。只有同时满足该要求，表~\ref{tab:load-dimensions} 中按未归一化传递函数计算的匹配族维数才适用。当电阻变化时，匹配的响应核给出式~\eqref{eq:weighted-average}，其中的权重仍取决于激励。

\subsection{负载改变后的共模瞬态泄漏}
\label{sec:cm-pulse}

对于在 $g_0$ 下匹配、在 $g$ 下观测的支路，记 $\Delta B=B_P-B_N$，由式~\eqref{eq:general-load-shift} 得到 $F=G_{DC}=-(g-g_0)\Delta B/(D_PD_N)$。由于 $F(0)=0$ 且极点稳定，幅度为 $V_c$ 的共模阶跃产生衰减的差分响应。以 $p_k$ 表示采用物理频率单位、约分后仍保留的互异极点，以 $a_k$ 表示 $F(s)/s$ 的留数，则有
\begin{equation}\label{eq:cm-pulse}
y(t)=V_c\sum_k a_ke^{p_kt},\qquad
\int_0^\infty y(t)^2\,\mathrm dt
=V_c^2\sum_{k,j}\frac{a_ka_j}{-(p_k+p_j)}.
\end{equation}
若改以 $r_k$ 表示 $F$ 的留数，则 $a_k=r_k/p_k$。冲激响应能量的求和项采用 $r_kr_j$，与上述阶跃响应能量是不同的量。此外，$\sup|y|\le |V_c|\|h_{DC}\|_1$，其中 $h_{DC}=\mathcal L^{-1}F$。存在重极点时，需要使用相应的多项式与指数函数乘积项。

表~\ref{tab:branches} 中的 P/T 支路在 $g_0=0.1$ 下匹配，在 $g=0.2$ 下观测时，P 减去 T 的传递函数为
\begin{equation}\label{eq:dynamic-pt}
F(x)=\frac{0.1x}{(x^2+3.2x+1.4)(x^2+3.3x+1.4)}.
\end{equation}
其无量纲极点为 $-2.8,-2.677033,-0.522967,-0.5$。取 $T=10~\mu\mathrm s$，单位阶跃响应在 $28.20333~\mu\mathrm s$ 达到峰值 $0.009197601$（图~\ref{fig:dynamic-pulse}）；约分后保留的最慢极点给出 $20~\mu\mathrm s$ 的渐近时间常数。阶跃响应能量为 $3.17497\times10^{-9}V_c^2~\mathrm s$；要确定有限的稳定时间，还需指定容差。




\subsection{瞬态验证}
\label{sec:dynamic-verif}

指数激励的评估窗口为 $2$--$3$~ms，最大步长为 $0.1~\mu\mathrm s$；正弦激励的窗口为 $20$--$40$~ms，最大步长为 $1~\mu\mathrm s$。细化后的最大步长分别为 $0.05~\mu\mathrm s$ 和 $0.5~\mu\mathrm s$。

六个基准原生瞬态算例及两个步长细化算例检验了三级和四级归一化元件值均为 1 的梯形网络，采用 $R_*=1~\mathrm{k\Omega}$、$C_*=10~\mathrm{nF}$ 及 $10~\mathrm{k\Omega}$ 端接负载。四个匹配通道算例采用 $R=10~\mathrm{k\Omega}\,e^{200t}$，激励分别为恒定电流和与 $e^{1000t}$ 成正比的电流。在 $2.5~\mathrm{ms}$ 时，改变激励使匹配通道读数在 $N=3$ 时变化 $5.11377~\Omega$，在 $N=4$ 时变化 $11.94461~\Omega$，与式~\eqref{eq:dynamic-exp} 一致。

两个恒定电阻算例采用 $R=100~\mathrm{k\Omega}$、$I=8~\mu\mathrm A[1+0.2\sin(2\pi100t)]$，并将参考通道末级电容增大 $10\%$。两者分别得到 $\Delta\tau=-2.30769~\mu\mathrm s$ 和 $-2.85714~\mu\mathrm s$。保留符号的一阶预测与精确响应之间的相对均方根偏差分别为 $5.59\%$ 和 $8.63\%$。所有算例中，初始过渡结束后，仿真比值读数与精确值的最大偏差小于 $1.23\times10^{-6}~\Omega$。两个 $N=4$ 的时间步长加密算例使读数变化小于 $1.4\times10^{-6}~\Omega$。

\begin{figure}[htbp]
\centering
\includegraphics[width=0.67\linewidth]{figures/Fig_08_dynamic_weighting.pdf}
\caption{匹配三阶、四阶梯形支路中电阻比值对电流历程的依赖。纵轴为同一电阻 $R(t)=10\,\mathrm{k\Omega}\exp(200t)$ 在两种电流 $I_\alpha(t)=1\,\mathrm{\mu A}\exp(\alpha t)$ 下的比值差 $\widehat R_{\alpha=1000}(t)-\widehat R_{\alpha=0}(t)$，其中 $t$ 以秒计，$\alpha$ 的单位为 $\mathrm{s}^{-1}$。各串联电阻均为 1~k$\Omega$，各对地电容均为 10~nF，两条负载均为 10~k$\Omega$。实线为指数输入的解析响应，空心标记为原生 LTspice 比值插值后的差值。}
\label{fig:dynamic-weighting}
\end{figure}

\begin{figure}[htbp]
\centering
\includegraphics[width=0.67\linewidth]{figures/Fig_09_dynamic_slew.pdf}
\caption{四阶梯形支路对在恒定电阻 $R=100$~k$\Omega$ 和电流 $I(t)=8\,\mathrm{\mu A}[1+0.2\sin(2\pi\,100t)]$ 下的电阻比值误差，其中 $t$ 以秒计。两条支路均采用 1~k$\Omega$ 电阻、10~nF 电容和 10~k$\Omega$ 负载，参考支路末级对地电容增至 11~nF。黑线为精确稳态相量比值，红线为一阶矩预测 $-R\Delta\tau\dot I/I$，空心圆为原生 LTspice 数据。其中 $\Delta\tau=\tau_P-\tau_N=-2.85714~\mathrm{\mu s}$。}
\label{fig:dynamic-slew}
\end{figure}

\begin{figure}[htbp]
\centering
\includegraphics[width=0.67\linewidth]{figures/Fig_10_dynamic_pulse.pdf}
\caption{P/T 支路对在 5~k$\Omega$ 对地负载下的正向共模转换脉冲 $v_{oP}-v_{oT}$。共模输入从 1.0005~ms 开始，在 10~ns 内由 0 升至 1~V；横轴为相对该上升沿起点的时间。实线为有限斜坡输入的解析响应，空心圆为逐点读取原生 LTspice 文件得到的输出差值。在 10~k$\Omega$ 负载下匹配的设计，负载改为 5~k$\Omega$ 后产生瞬态输出差。}
\label{fig:dynamic-pulse}
\end{figure}

\FloatBarrier

\section{仿真结果与输出耦合}

\subsection{端接条件对等效关系的影响}

表~\ref{tab:branches} 汇总三组支路构造。各组等式与不等式关系均在扫频前通过代数运算确定。

\begin{table}[htbp]
\centering
\caption{归一化元件值与支路系数。}
\label{tab:branches}
\begin{tabular}{lrrrrrrrr}
\toprule
支路 & $r_1$ & $r_2$ & $q_1$ & $q_2$ & $a$ & $b$ & $c$ & $d$ \\
\midrule
P & 1 & 1 & 1 & 1 & 3 & 1 & 2 & 1 \\
Q & 1 & 4 & $1/2$ & $1/2$ & 3 & 1 & 5 & 2 \\
T & $18/19$ & $20/19$ & $361/180$ & $1/2$ & 2.9 & 1 & 2 & 2 \\
\bottomrule
\end{tabular}
\end{table}

图~\ref{fig:termination-transfer} 比较 P/Q 的开路和带载响应。图~\ref{fig:output-impedance} 给出区分这两条开路等效支路的输出阻抗。图~\ref{fig:load-specific} 比较两种负载下 T 相对于 P 的响应。



仿真复现了 P/Q 的开路共模抵消，以及 P/T 在 \(R_L=10\,\mathrm{k\Omega}\) 下的共模抵消。将 \(R_L\) 改为 5 k\(\Omega\) 后，P/T 响应不再相等。在全部 37 个算例中，LTspice 输出与独立节点方程解之间的复节点电压绝对差最大值为 \(1.14\times10^{-15}\) V。

\begin{figure}[htbp]
\centering
\includegraphics[width=0.67\linewidth]{figures/Fig_03_termination_transfer.pdf}
\caption{P、Q 在开路和 10~k$\Omega$ 对地负载下的幅频响应。实线为解析传递函数，空心圆为已有 LTspice 仿真数据。两条开路响应重合。归一化基准为 $R_*=1$~k$\Omega$、$C_*=10$~nF。}
\label{fig:termination-transfer}
\end{figure}

\begin{figure}[htbp]
\centering
\includegraphics[width=0.67\linewidth]{figures/Fig_04_output_impedance.pdf}
\caption{理想输入电压源置零后，P、Q、T 的开路输出阻抗幅值。曲线由四系数表达式计算。P 与 Q 的开路电压传递函数相同，输出阻抗不同。}
\label{fig:output-impedance}
\end{figure}

\begin{figure}[htbp]
\centering
\includegraphics[width=0.67\linewidth]{figures/Fig_05_load_specific_match.pdf}
\caption{两种共用对地负载下 T 相对于 P 的幅值响应。纵轴为 $20\log_{10}|H_T/H_P|$，零值表示幅值相等。实线为解析结果，空心圆为已有 LTspice 仿真数据。10~k$\Omega$ 下完整传递函数的相等性由正文代数推导确定。}
\label{fig:load-specific}
\end{figure}


\subsection{跨接电容同时作用于两种差分输出增益}

保持理想电压源和指定并联负载不变，在两个输出节点之间接入电容 \(C_d\)。以 \(Z'_P\) 和 \(Z'_N\) 分别表示各支路输出阻抗与本侧负载并联后的阻抗，以 \(H_P\) 和 \(H_N\) 表示接入 \(C_d\) 前的带载传递函数。联立两个输出节点方程可得

\begin{equation}\label{eq:12}
\begin{aligned}
G_{DC}&=\frac{H_P-H_N}{D_d},\\
G_{DD}&=\frac{(H_P+H_N)/2}{D_d},\\
D_d&=1+sC_d(Z'_P+Z'_N).
\end{aligned}
\end{equation}

在 \(G_{DD}/G_{DC}\) 有定义的频点，共同因子 \(D_d\) 从比值中约去。跨接电容以同一因子改变两种差分输出增益，因而保持 \(G_{DD}/G_{DC}\) 不变。若原带载支路传递函数相等，接入跨接电容后仍有 \(G_{DC}\) = 0。

对于 \(R_L=10\,\mathrm{k\Omega}\)、\(C_d=820\,\mathrm{pF}\) 的 P/Q 网络，在 145 个频点上，相对于未接跨接电容时复增益比的最大相对差为 \(5.14\times10^{-12}\)。该结果与式~\eqref{eq:12} 中共同因子的约去一致。


图~\ref{fig:bridge-gains} 展示两个独立增益的变化，图~\ref{fig:bridge-cmrr} 比较其共模抑制比。

\begin{figure}[htbp]
\centering
\includegraphics[width=0.67\linewidth]{figures/Fig_06_bridge_gains.pdf}
\caption{P/Q 支路对在有、无跨接电容时的解析差分输出增益。DM 表示单位差模输入；CM$\rightarrow$DM 表示单位共模输入产生的差分输出。空心标记用于识别 820~pF 曲线，其数值同样来自解析计算。两侧对地负载均为 10~k$\Omega$。}
\label{fig:bridge-gains}
\end{figure}


\begin{figure}[htbp]
\centering
\includegraphics[width=0.67\linewidth]{figures/Fig_07_bridge_cmrr.pdf}
\caption{负载条件下 P/Q 支路对的共模抑制比。$C_d=0$ 和 820~pF 的解析曲线重合。空心圆由 $C_d=820$~pF 时已有共模和差模 LTspice 输出差值计算。两侧对地负载均为 10~k$\Omega$；CMRR 为 $20\log_{10}|G_{DD}/G_{DC}|$。图~\ref{fig:bridge-gains} 同时给出绝对增益，图~\ref{fig:terminal-sweep} 展示独立支路的端接扫描。}
\label{fig:bridge-cmrr}
\end{figure}


\FloatBarrier

\section{讨论}

负载规范决定哪些元件组合满足共模抵消要求。开路共模抵消固定两个系数，指定正电阻负载固定三个系数组合，两个不同正电阻负载则固定全部四个系数。元件恢复不等式进一步将这些系数要求转化为可直接检验的正值元件设计域。

P/Q 与 P/T 的区别给出了直接的设计步骤。对于开路目标，根据所需电压传递函数确定 a、b，再在式~\eqref{eq:8} 内选择 c、d。对于单一指定电阻负载，确定 b、c 和 a + gd，并利用剩余自由度构造支路。对于覆盖两个不同电阻负载的抵消要求，该拓扑必须采用相等的对应元件值。

\citet{casas1996basics} 在另一种电路配置中通过浮置耦合实现共模抵消。本文的对地梯形网络保留了流向地的电流路径，式~\eqref{eq:6} 显式给出了端接负载的作用。\citet{hornero2021cm} 通过共模响应描述了随频率变化的失配现象。本文的系数表示将完整的二阶支路响应与元件值、输出阻抗联系起来。

元件恢复结论要求元件值严格为正。当 \(q_1\) = 0 时，两个串联电阻合并为一个电阻；元件组合 (1,2,0,3) 和 (3/2,3/2,0,3) 得到相同的系数 (9,0,3,0)，式~\eqref{eq:7} 此时出现奇异。因此，双负载唯一性结论适用于第~\ref{sec:model-plane} 节定义的全正值二阶梯形网络。

对于导纳为 \(Y_L(s)\) 的频率相关线性负载，支路传递函数为 \(H_0/[1+Z_0Y_L(s)]\)。设计面向此类负载时，可先代入其导纳，再比较两个支路的响应。式~\eqref{eq:12} 则描述两个输出节点之间的耦合情形。


\subsection{设计与验证步骤}

滤波器规范可分四步转化为元件约束。首先，确定信号源平面、输出测量平面，以及各输出端所见的负载导纳。其次，选择该端接下所需的分母系数，并在正值元件域内选定剩余自由系数。再次，按式~\eqref{eq:7} 恢复元件值，并逐项检验式~\eqref{eq:8} 的不等式。最后，分别扰动元件值和两侧负载导纳，计算相应的 $G_{DC}$ 与 $G_{DD}$。精确等式确定标称设计；式~\eqref{eq:sensitivity} 和式~\eqref{eq:load-asymmetry} 则指出需要控制容差的组合。

双负载结论同时提供了一种验证思路。在两个已知电阻负载下测量完整支路响应，可约束四个系数。实验中需根据复频率响应数据估计这些系数，并给出不确定度区间。精确定理确定可辨识的模型结构；两个相近元件组合能够被区分到何种程度，则取决于测量不确定度与所选频点。

\subsection{待完成的实验验证}

\begin{quote}
\textbf{待补实验数据。} 当前结果包括精确代数推导与电路仿真。下一版预留 P/Q、P/T 在开路、10~k$\Omega$ 和 5~k$\Omega$ 端接下的实物电路比较。元件实测值、信号源与探头的复阻抗、频率响应、不确定度预算和重复性记录尚未采集。
\end{quote}

实验将采用相同的四元件梯形网络，以及表~\ref{tab:branches} 的元件比例。先分别测量各支路响应，再对支路对施加共模和差模激励。记录电阻、电容实际值后，可分别比较标称元件预测与实测元件预测。负载和探头导纳按输出测量平面接入带载模型。对于接近仪器噪声底的共模泄漏，将报告带界限的测量结果及相应检出限。

安装跨接电容后将进行第二轮采集，比较两种复增益及其比值。将信号源和探头配置表征为两个独立的线性支路等效网络，利用其带载传递函数和输出阻抗检验式~\eqref{eq:12} 的共同因子预测。记录的响应将与实测元件模型预测进行比较，并结合两者的不确定度判断符合程度。

\section{结论}

本文建立了对地差分电阻电容梯形网络随端接条件变化的共模抵消条件：先在二阶完成全细节推导，再以链式矩阵递推推广到任意级数。参考支路固定时，开路抵消对应双参数匹配族，单一正电阻端接对应单参数匹配族；在两个不同正电阻端接下均具有相同完整传递函数，则要求对应元件值相同。正值元件恢复条件和显式支路构造使各项结论均可直接检验。37 个二阶与 84 个一般阶交流仿真算例复现了解析预测。输出跨接电容在两种差分输出增益中引入相同负载因子，从而保持共模抑制比不变。对时变测量，负载传函相等被证明恰是比值读出对任意激励电流动态精确免疫的条件；一阶交叉项与共模瞬态泄漏脉冲量化了端接失配的动态代价。



\FloatBarrier
\bibliographystyle{plainnat}
\bibliography{references}
\end{document}
